%% Citas verificadas contra el PDF. Sin marcadores [VERIFY] pendientes.

\section{VALIDACIÓN}

Esta sección describe el protocolo con el que se someterá a prueba la arquitectura completa.
Como meta piloto inicial se fija que la extensión resuelva al menos el 70\,\% de las
características en fichas de plataformas grandes y al menos el 40\,\% en vendedores
independientes. Son valores de referencia para orientar la primera medición, no umbrales de
éxito acreditados por evidencia previa, y se revisarán con los datos que arroje la Fase~1. No
se fijan aquí metas sobre la proporción de advertencias infundadas ni sobre el efecto en la
decisión de compra: ambas exigen una referencia externa sobre la legitimidad de cada ficha de
la que este protocolo todavía no dispone.

\subsection{Verificación preliminar del vector de características}

Para verificar que el vector de ocho características de la Tabla~\ref{tab:senales} contiene
información discriminante, se construyó un conjunto sintético de 10\,000 registros con tres
clases equilibradas. El conjunto se construyó ante la ausencia de un conjunto etiquetado
público de fichas de vendedores peruanos: sus variables y rangos se derivaron de los factores
de riesgo identificados en la Sección~II, y las etiquetas se asignaron mediante reglas
deterministas sobre esas mismas variables. Declarar esta procedencia es indispensable para
interpretar lo que sigue. Sus siete entradas corresponden a siete de las ocho de la
Tabla~\ref{tab:senales}. La veracidad de dominio, que este trabajo introduce, es la única que
no forma parte de ese vector. Se entrenaron cinco algoritmos de familias distintas sobre la
misma partición 80/20 con semilla fija, y se validaron además mediante validación cruzada
estratificada de cinco pliegues sobre el conjunto completo. Los resultados muestran que los
cinco algoritmos clasifican correctamente los 2\,000 registros de prueba, sin desviación
entre pliegues.

Los experimentos se ejecutaron sobre Ubuntu 24.04.3 LTS con Python 3.11 y scikit-learn, en un
procesador Intel Xeon E3-1240 v5 y sin aceleración gráfica. Las máquinas de soporte vectorial y
el clasificador de vecinos más cercanos se encadenaron a un normalizador de media cero y
varianza unitaria; los demás hiperparámetros se mantuvieron por defecto, con semilla fija,
porque el objetivo no era optimizar el desempeño sino comprobar si existía diferencia alguna
entre familias de algoritmos.

Los cinco algoritmos recuperan la regla generadora: que un clasificador de vecinos más
cercanos iguale a un ensamble de gradiente potenciado, y que ambos igualen a un árbol
individual, indica que la frontera de decisión es recuperable por cualquier método, de modo
que el problema es separable por construcción. Un árbol de decisión de profundidad cinco
reproduce la regla de etiquetado en su totalidad, y la importancia de las características se
concentra en dos: la puntuación del producto, con 0.51, y el volumen de ventas, con 0.46,
mientras que tres de las siete no contribuyen en absoluto.

El ejercicio vale como \textbf{verificación del canal de procesamiento y diagnóstico del
conjunto}, no como medida de eficacia. De ello se siguen dos consecuencias. La primera es que
estos datos \textbf{no permiten discriminar entre algoritmos}: la elección del bosque
aleatorio descrita en la Sección~III responde a requisitos de interpretabilidad y costo, no a
una comparación empírica que este conjunto no puede arbitrar. La segunda, más relevante, es
que las métricas obtenidas sobre él no estiman la calidad del nivel estimado de riesgo, sino
la fidelidad con que un modelo recupera la regla que generó sus propias etiquetas. La validez
del enfoque depende de que las características sean legibles en fichas reales y de que las
etiquetas provengan de una fuente independiente del vector de entrada, según describe el
protocolo de la siguiente subsección.


\subsection{Diseño experimental de la validación}

La validación se ejecutará en el Perú, sobre las plataformas y los hábitos descritos en la
Sección~I. Esa elección delimita el alcance de las conclusiones empíricas, no el de la
arquitectura: las características de la Tabla~\ref{tab:senales} se derivan de elementos
presentes en cualquier ficha de producto y no de particularidades del mercado peruano. De ahí
se sigue una hipótesis de transferibilidad, no una propiedad demostrada: que la propuesta
opere en otros mercados depende de que esas características resulten observables en ellos, lo
que exige medir su cobertura en cada caso.

La validación se organiza en dos fases con objetivos separables.

\emph{Fase 1: observabilidad de las características.} Antes de evaluar la utilidad del
veredicto es necesario establecer qué fracción de las características de la
Tabla~\ref{tab:senales} resulta efectivamente legible en fichas reales. Sobre una muestra de
fichas de producto de las plataformas de mayor uso en el Perú se registrará, para cada
característica, si pudo resolverse o si quedó indeterminada, distinguiendo entre plataformas
grandes y vendedores independientes. El desenlace relevante no es un porcentaje agregado sino
su distribución: una arquitectura que solo resolviera características en las plataformas que ya
ofrecen garantías carecería de utilidad precisamente en el segmento donde el riesgo se
concentra.

\emph{Fase 2: comportamiento de compra con la herramienta.} Un diseño con dos grupos, con y
sin la extensión activa, enfrentará a los participantes al mismo conjunto de fichas reales. La
asignación al grupo será aleatoria y el orden de presentación de las fichas se
contrabalanceará entre participantes, para que el aprendizaje sobre la tarea no se confunda
con el efecto de la herramienta. La población objetivo son compradores digitales peruanos de
entre 18 y 60 años. La muestra prevista de cincuenta participantes es una muestra piloto,
dimensionada por factibilidad: permite estimar la variabilidad de las medidas y detectar
problemas del protocolo, no alcanzar potencia estadística, que deberá calcularse a partir de
esa variabilidad antes del estudio definitivo. Todos otorgarán su consentimiento de forma
voluntaria antes de comenzar, y las sesiones se conducirán de manera presencial o en línea,
según su disponibilidad. Los registros de consulta descritos en la Sección~III se recogerán
disociados del participante que los originó, se conservarán únicamente durante el periodo de
análisis y se eliminarán al cierre del estudio; ningún dato que permita reconstruir el
historial de compra de una persona identificada forma parte de lo que el protocolo almacena.

Las medidas de esta fase se eligieron bajo una restricción que conviene declarar: el diseño no
dispone de una verdad de referencia sobre cada ficha. Establecer si un vendedor es fraudulento
exige evidencia sobre operaciones consumadas denuncias verificadas, retiros de la plataforma
que la ficha no contiene y que este protocolo no produce. Por eso no se miden aquí la tasa de
identificación correcta ni la tasa de falsos positivos sobre comercios legítimos: ambas
presuponen esa referencia, y construirla es un requisito previo del estudio de eficacia, no un
subproducto de este.

Se medirá, en cambio, lo que el diseño sí puede observar. Del comportamiento: la decisión de
compra sobre cada ficha en uno y otro grupo, la proporción de decisiones que el participante
revierte tras leer el veredicto y el tiempo hasta decidir. Del sistema: la cobertura de
características alcanzada en cada ficha y la proporción de fichas en que la regla se abstiene,
magnitud que gobierna la utilidad del instrumento, pues una herramienta que se abstiene casi
siempre no llega a informar decisión alguna. La percepción del participante se recoge mediante
un cuestionario estructurado de quince ítems con respuesta en escala de Likert de cinco
puntos, organizado en las cinco dimensiones de la Tabla~\ref{tab:instrumento}. Solo el bloque
de usabilidad emplea un instrumento validado, la Escala de Usabilidad del Sistema
\cite{ref23}; las cuatro dimensiones restantes constituyen un instrumento propio, sujeto a
validación en el propio estudio piloto, del que se reportará la consistencia interna antes de
interpretar sus resultados. La dimensión de influencia en la decisión es autorreportada y no
se emplea como indicador de eficacia: que una advertencia influya no implica que la decisión
resultante sea mejor, y establecerlo exige la referencia externa que esta fase no tiene.

\begin{table}[t]
\caption{Dimensiones del instrumento de validación con usuarios}
\label{tab:instrumento}
\centering
\footnotesize
\setlength{\tabcolsep}{3.5pt}
\renewcommand{\arraystretch}{1.15}
\begin{tabular}{@{}p{2.35cm}p{5.3cm}@{}}
\toprule
\textbf{Dimensión} & \textbf{Qué establece} \\
\midrule
Usabilidad & Si el veredicto y su explicación se interpretan sin instrucción previa (bloque SUS) \\
Confianza percibida & Si el participante considera fundada la advertencia que recibe \\
Riesgo percibido & Si su estimación del riesgo de la ficha varía al disponer de la herramienta \\
Intención de uso y recomendación & Si declararía seguir usándola fuera del entorno de prueba \\
Influencia en la decisión & Si la advertencia modificó lo que habría hecho sin ella \\
\bottomrule
\end{tabular}
\end{table}

La asimetría de costos que motiva la regla de abstención de la Sección~III se mantiene como
principio de diseño y no como resultado medible en esta fase. Un instrumento que advierte con
excesiva frecuencia pierde credibilidad y su costo recae sobre comercios honestos; cuantificar
ese costo exige la verdad de referencia que el estudio de eficacia deberá construir. Por eso
el veredicto de riesgo y el de evidencia insuficiente se mantienen separados: son resultados
distintos y su evaluación no puede confundirse.

\subsection{Amenazas a la validez}

Cuatro limitaciones condicionan lo que este diseño podrá concluir, y conviene explicitarlas
antes de su ejecución.

La primera afecta a la capa cliente. La extracción depende de selectores sobre el DOM, y se
ha advertido que las técnicas apoyadas en selectores predefinidos son sensibles a cambios de
maquetación \cite{ref24}. Una plataforma que rediseñe su ficha reducirá la cobertura de
características sin previo aviso, por lo que la Fase~1 debe repetirse periódicamente y no
tratarse como una medición única.

La segunda afecta al clasificador de riesgo. Sus etiquetas de entrenamiento se derivan
mediante reglas deterministas aplicadas a las mismas características que el modelo recibe como
entrada, de modo que el desempeño medido sobre ese conjunto refleja la recuperación de esa
regla y no el comportamiento frente a vendedores fraudulentos reales. La mitigación es previa
a cualquier medición de eficacia y condiciona el orden del trabajo futuro: construir un
conjunto cuyas etiquetas provengan de una fuente independiente del vector de características,
y hasta entonces no reportar métricas de desempeño del clasificador.

La tercera es de validez ecológica. Un experimento en sesión mide la capacidad de decisión en
condiciones de atención deliberada, mientras que la evidencia disponible sugiere que la
victimización ocurre justamente cuando esa atención falla \cite{ref27}. El diseño puede
establecer si la herramienta ayuda a quien la mira, no si protege a quien compra distraído.
Mitigarlo exige una segunda medición en uso real, con la extensión instalada durante semanas y
sobre las compras que el participante hubiera hecho de todos modos, que este protocolo no
cubre.

La cuarta es la más general: la arquitectura supone que el atacante no adapta su
comportamiento a las características observadas, suposición insostenible a mediano plazo
porque los defraudadores ajustan sus estrategias frente a los mecanismos de detección
\cite{ref28}. Como mitigación parcial, el registro de consultas descrito en la Sección~III
hace observable esa deriva y permite reentrenar cuando la distribución de las características
cambie; ninguna medición puntual, en cambio, conserva su validez indefinidamente.
